% sec-06: Document 06, the state regulation under Title 22.
% STAGE 3.  The ten-row staffing table moved to mvltable, because a ten-row
% table with three-line cells is taller than the space a section heading and two
% paragraphs leave on the page it opens.  Two column widths retuned.
\docpart{Title 22 Chapter 15}{California Code of Regulations, Title 22,
Division 5, Chapter 15 - PDAC LLM Oncology Clinical Trial Site Regulations}

\section{Scope and Authority}

This chapter is adopted under the authority of document 01 of this package and
the licensing provisions of the Health and Safety Code \cite{cahsc,ccr22}. It
governs the authorization, survey, and enforcement of a pancreatic oncology
clinical trial site incorporating an LLM-directed clinical workflow.

This chapter does not govern a general acute care hospital, a clinic licensed
under other chapters of this division, or a research laboratory that enrolls no
human participants. Where a facility holds another license, this chapter applies
only to the trial activity described in document 01 §2 and does not displace the
requirements of that license.

\section{Application for Authorization}

An application shall contain each of the following. Each item is an artifact
that exists elsewhere in this package, so that an applicant is never asked for a
document that has not been defined.

\begin{mvtable}
\begin{tabularx}{\textwidth}{>{\raggedright\arraybackslash}p{0.9cm} >{\raggedright\arraybackslash}p{4.4cm} >{\raggedright\arraybackslash}p{2.4cm} >{\raggedright\arraybackslash}X}
\headrow \thead{Item} & \thead{Content} & \thead{Source} & \thead{Form in which it is filed}\\
\midrule
(a) & Staffing plan and delegation log & Document 14 & Roster, qualifications, signed delegation entries \\
(b) & Premises plan with zones and envelopes & Document 09 & Floor plan with zone boundaries and envelope markings \\
(c) & Building commissioning record & Document 08 §6 & Commissioning report and certificate of occupancy \\
(d) & Emergency preparedness plan & Document 11 & Plan, drill schedule, mutual aid letter \\
(e) & Model version register & Document 03 §3 & Register extract with content hashes \\
(f) & Robot instance register & §4 of this chapter & Serial numbers, envelope, last interlock test \\
(g) & Protocol and regulatory status & Document 13 §2 & Protocol, institutional review board approval, safe-to-proceed record \\
(h) & Standard operating procedure index & Document 12 §3 & Index with owners and review dates \\
(i) & Parking and transportation plan & Document 10 & Stall schedule and the transportation demand plan \\
\end{tabularx}
\end{mvtable}
\tabcapl{tab:t22app}{Nine application items. A regulation whose application asks
for documents that do not exist cannot be complied with, so every row names the
document in this package that produces the artifact.}

\section{Staffing and Qualification Requirements}

\begin{mvltable}
\begin{xltabular}{\textwidth}{>{\raggedright\arraybackslash}p{4.9cm} >{\raggedright\arraybackslash}p{1.8cm} >{\raggedright\arraybackslash}X}
\headrow \thead{Position} & \thead{Minimum} & \thead{Qualification the department will verify}\\
\midrule
\endfirsthead
\headrow \thead{Position} & \thead{Minimum} & \thead{Qualification the department will verify}\\
\midrule
\endhead
Site principal investigator & 0.10 FTE & California medical license; hepatopancreatobiliary surgical practice; 40 hours investigator training \\
Sub-investigator, medical oncology & 0.10 FTE & California medical license; gastrointestinal medical oncology practice; 40 hours investigator training \\
Director of clinical operations & 0.40 FTE & Five years of trial operations; 80 hours coordinator training \\
Lead clinical research coordinator & 1.00 FTE & Coordinator certification; 80 hours training; consent competency assessment \\
Regulatory affairs and quality manager & 0.45 FTE & Three years of regulatory submissions; 80 hours training \\
Investigational drug pharmacist & 0.20 FTE & California pharmacist license; sterile compounding competency under USP 797 and 800 \\
Site safety officer & 0.55 FTE & 160 hours safety training; robot platform certification; interlock test competency \\
Model governance lead & 0.40 FTE & 120 hours training; verification and validation competency; audit trail administration \\
Data manager and biostatistician & 0.30 FTE & Statistical qualification; database lock competency \\
Research nurse and participant navigator & 0.25 FTE & California registered nurse license; oncology experience; 80 hours training \\
\end{xltabular}
\end{mvltable}
\tabcapl{tab:t22staff}{Staffing minimums by position. The full-time equivalent
column is identical to the roster in document 14 §1 and the training hours are
identical to document 12 §5; the department surveys against those two documents
rather than against a separate standard.}

The chief executive and sponsor representative is not listed above, because that
role is not a clinical position and the department does not verify a
qualification for it. Its exclusion is deliberate and is explained in document
14 §3.

\section{Robot and Model Registration}

\begin{enumerate}
\item A site shall maintain a robot instance register with one row per physical
  robot: manufacturer, model, serial number, suite assignment, envelope
  identifier, date of last interlock test, and date of last envelope survey. The
  site described in this package registers 14 instances across 8 types.

\item A site shall maintain a model version register under document 03 §3, with
  one row per model version placed into clinical use.

\item An amendment to either register shall be filed with the department within
  10 business days of the event that triggers it. The triggering events are:
  addition or removal of a robot instance; relocation of an instance between
  suites; a change to an envelope; a failed interlock test; and the placement of
  a new model version into clinical use.

\item A robot instance that has failed an interlock test shall not be used in a
  study procedure until it has passed a retest and the register has been
  amended.
\end{enumerate}

\section{Survey, Deficiency Classes, and Penalties}

\begin{enumerate}
\item The department shall survey each authorized site at least annually, and
  may survey without notice at any time during the hours of operation stated in
  document 05 §5.

\item Deficiencies are classified as Class A, Class B, or Class C with the
  meanings, correction periods, and civil penalty ranges established in document
  01 §5. Those definitions are not restated here and are not varied.

\item The department shall issue a statement of deficiencies within 15 business
  days of the survey exit conference. The site shall file its correction plan
  within the period the class requires.

\item An appeal is taken under document 01 §5(c). A Class A suspension is not
  stayed by an appeal.
\end{enumerate}

\section{Records and Reporting}

\begin{enumerate}
\item A site shall file an annual report stating: participants screened,
  enrolled, and treated; every reportable robotic event under document 02 §2;
  every model fault under document 11 §5; every deficiency cited and its
  correction; and the current registers.

\item Records shall be retained for six years from database lock, which is the
  period fixed by document 03 §5. Where this chapter, federal law, and the
  facility record law of this state prescribe different periods, the longest
  governs.

\item A site shall notify the department within 24 hours of any event that would
  constitute a Class A deficiency, whether or not the department has cited it.
  Self-reporting within 24 hours is a mitigating factor in the civil penalty
  determination and shall be stated as such in the department's order.
\end{enumerate}
